\subsubsection{Proof of the remaining eleven exclusions}
\label{sec:m9-completion}

\begin{theorem}\label{thm:m9-completion}
None of the eleven rings
\[
 \RR{9}{529},\ \RR{9}{1060},\ \RR{9}{2455},\ \RR{9}{3748},\
 \RR{9}{7018},\ \RR{9}{7741},\ \RR{9}{7742},\
 \RR{9}{8099},\ \RR{9}{8101},\ \RR{9}{8200},\ \RR{9}{8785}
\]
admits a fusion categorification over an algebraically closed field of
characteristic zero. Consequently all $8\,809$ exact-multiplicity-nine
rings are excluded, and none contributes a positive model.
\end{theorem}

We prove the remaining cases by three common reductions: arithmetic of
signed dimension fields, a small Casimir-deficit argument, and central
induction. All finite domains below are derived from these reductions.
The complete tensors, integer witnesses and independently executable
verifiers are in \filename{research/exclusions/m9_completion/}.
No arbitrary conductor cutoff or unfinished associator search is used.

\paragraph{Signed necessary data and two useful refinements.}
For a self-dual input ring, sphericalization of an arbitrary hypothetical
category gives the commutative signed algebra of
Lemma~\ref{lem:dimension-field-containment}. Indeed, the two lifts of a
self-dual simple have opposite nonzero dimensions, whereas duality
preserves dimension, so both lifts are self-dual. Its coefficients satisfy
\[
 t_{ijk}=t_{\pi(i,j,k)},\qquad |t_{ijk}|\le N_{ijk},\qquad
 t_{ijk}\equiv N_{ijk}\pmod2.
\]
Write $K$ for its dimension-character field, $\Delta$ for its dimension
codegree and $T_0=\sum N_{ijk}^2$. The symmetric integral regular matrices
give a totally real semisimple algebra. Cyclotomicity makes $K$ abelian.
Every original or signed codegree $f$ belongs to $K$, and
$\Delta/f\in\mathcal O_K$: apply central induction to sphericalization,
whose corresponding codegrees and global dimension are $2f$ and
$2\Delta$. The six signed codegrees exceed one, have reciprocal sum one,
are $d$-numbers, and have total at most $T_0$. The conjugates of $\Delta$
contribute $\operatorname{Tr}_K(\Delta)$ to this total, even when
$\Delta$ itself generates a proper subfield of $K$.
These statements are the earlier signed-field and trace-ideal criteria,
using \cite[Section~5.1 and Theorem~2.13]{Ostrik3} and
\cite[Theorem~1.2 and Corollary~1.4]{Ocode}.

\begin{lemma}[Indicator normalization]\label{lem:m9-indicator-gauge}
For the self-dual signed algebra above, the lifts can be chosen so that
\[
 N_{jji}-t_{jji}\equiv0\pmod4\qquad\text{for all }i,j.
\]
No further independent sign normalization is implicit in this choice.
\end{lemma}
\begin{proof}
The Fuchs--Runkel--Schweigert theorem, in the pivotal formulation
\cite[Theorem~3.9]{LPRV2025}, states that a self-dual simple of indicator
$-1$ occurs with even multiplicity in $Y^*\otimes Y$; neither braiding
nor simplicity of $Y$ is required. The two sphericalization lifts differ
by the sign of their pivotal map. The transpose operator
$\kappa\mapsto\kappa^*j$ on the one-dimensional self-duality space
therefore changes sign, so their indicators are opposite. Choose the
indicator-$+1$ lift of each simple. The negative lift of $i$ has
multiplicity $(N_{jji}-t_{jji})/2$ in the square of the chosen lift of $j$,
and that multiplicity is even. Apply the theorem inside sphericalization,
without requiring a pivotal structure on the original category.
\end{proof}

\begin{lemma}[Diagonal associativity and saturation]
\label{lem:m9-diagonal}
Every self-dual signed table as above satisfies
\[
 \sum_k |t_{iik}t_{jjk}|\ \geq\ \sum_k t_{ijk}^2.
 \tag{$\dagger$}
\]
If $|t|=N$ and every nonunit $N_{iik}$ is positive, a sign gauge makes
all repeated-index triples positive. Only the ten triples with three
distinct nonunit indices then remain to be tested.
\end{lemma}
\begin{proof}
The coefficient of $1$ in $(i^2)(j^2)=(ij)^2$ gives
$\sum_k t_{iik}t_{jjk}=\sum_k t_{ijk}^2$, proving the inequality.
In the saturated case the unsigned original table gives the same sum
with each summand replaced by its maximum $N_{iik}N_{jjk}$. Equality of
the sums forces equality termwise. Thus for fixed nonunit $k$ the signs
of $t_{iik}$ are independent of $i$. Gauge the basis element $k$ by this
common sign. The transformed $t_{iik}$ is positive, since the gauge on
$i$ occurs twice. This includes $i=k$.
\end{proof}

\paragraph{Three small-deficit cases.}
For $\RR{9}{7741}$ and $\RR{9}{7742}$ the trace-ideal lower bound is
$6705$, and the original Casimir traces are $6813$ and $6805$.
For an unordered nonunit triple $q$ of weight $w_q$ (its number of
ordered permutations), enumerate
\[
 a_q\in\{N_q,N_q-2,\ldots\}\cap\mathbb Z_{\geq0},\qquad
 \sum_qw_q(N_q^2-a_q^2)\le108\quad\text{or}\quad100.
\]
These are respectively $389$ and $241$ magnitude profiles. Every profile
except $a=N$ violates $(\dagger)$ for at least one pair $i,j$.
This test is independent of all coefficient signs. Both tensors have
$N_{iik}>0$, so Lemma~\ref{lem:m9-diagonal} leaves $2^{10}$ sign choices;
checking the complete associativity identities leaves only the positive
table. The resulting fusion-closed lifts give a spherical section.
The previously established complete spherical obstructions for these two
tensors therefore exclude arbitrary categorifications. Their exact
spherical proof trees are replayed as part of the new certificate.

For $\RR{9}{8200}$, put $E=\mathbb Q(\sqrt{401})$. The original
FP codegree is $D=10\sqrt{401}(20+\sqrt{401})$ and $T_0=8040$.
Since $D\in2\mathcal O_E$ and $\Delta/D$ is integral,
$\Delta\in2\mathcal O_K$. The characteristic polynomial
$\operatorname{Norm}_{K/\mathbb Q}(X-\Delta)$ divides the signed
Casimir polynomial and reduces to $X^{[K:\mathbb Q]}$ modulo two.
The original Casimir polynomial reduces to $X^2(X+1)^4$; hence
$[K:\mathbb Q]\le2$, and $E\subset K$ gives $K=E$.
The unique $\mathbb F_2$ character is $(1,0,0,1,0,0)$.
As two splits in $E$, reduction of every dimension embedding has this
character, so $\Delta=\sum_i d_i^2\equiv2\pmod{4\mathcal O_E}$.
Its trace is divisible by four. Combining this with the trace ideal
$4010\mathbb Z$ gives $\operatorname{Tr}_E(\Delta)=8020$.
The resulting deficit at most $20$ is smaller than the least possible
unsaturated loss $28$. Saturation and the same $2^{10}$ sign test give
a spherical section. Only the FP dimension orbit can occur; the quadratic
Gauss obstruction, Proposition~\ref{prop:quadratic-gauss-family}, then
requires at least $3200$ root contributions where only $770$ are available.
This excludes $\RR{9}{8200}$.

\paragraph{The two cyclic-quartic cases.}
For $i=8099,8101$ the original codegree polynomial is $F G_i$, where
\[
\begin{split}
 F&=X^2-4225X+21125,\\
 G_{8099}&=X^4-300X^3+5250X^2-22500X+28125,\\
 G_{8101}&=X^4-300X^3+26250X^2-562500X+703125.
\end{split}
\]
Each quartic is irreducible and generates
$E=\mathbb Q(r)$, $r^4-r^3-4r^2+4r+1=0$.
The integral basis is $1,r,r^2,r^3$, of discriminant $1125$.
The automorphism $r\mapsto-r^3-r^2+3r+2$ has order four, and
$1-6r+2r^3$ squares to five. Thus $E$ is cyclic quartic with unique
quadratic subfield $\mathbb Q(\sqrt5)$. Field containment gives $K=E$.
Intersecting the six original codegree ideals gives the column lattices
\[
 L_{8099}=65\begin{pmatrix}1&0&0&0\\0&1&0&0\\-8&-13&15&0\\-7&-4&0&15\end{pmatrix},
 \qquad
 L_{8101}=65\begin{pmatrix}1&0&0&0\\-2&5&0&0\\-12&-5&15&0\\1&-5&0&15\end{pmatrix}
 \subset L_{8099}.
\]
Consequently $\Delta=65\theta$, where
$\theta=a+br+cr^2+dr^3$ satisfies
\[
 c\equiv-8a-13b,\quad d\equiv-7a-4b\pmod{15},\qquad
 t:=\operatorname{Tr}(\theta)=4a+b+9c+d\in\{15,30,45,60\}.
\]
Here $T_0=4525$ and the trace ideal is $975\mathbb Z$.
Every conjugate of every trace-dual basis element lies between $-2$ and
$2$ (checked by exact Sturm counts), so total positivity gives
$|a|,|b|,|c|,|d|<2t$. This is an exhaustive finite box.
For $P_\theta=X^4+c_1X^3+c_2X^2+c_3X+c_4$, require
$c_4^j\mid c_j^4$ for $1\le j\le4$, the $d$-number conditions.
They remain necessary if $P_\theta$ is a power of a smaller minimal
polynomial. The two remaining signed codegrees $f,g$ form a rational
pair or a conjugate pair in $\mathbb Q(\sqrt5)$. In lowest terms put
\[
 \frac AB=1+\frac{c_3}{65c_4}=\frac1f+\frac1g>0.
\]
Then $f+g=Ak$, $fg=Bk$, with $k\ge1$ and $Ak\le4525-65t$.
Both roots of $X^2-AkX+Bk$ must exceed one. Its discriminant is a square
with integral roots, or five times a square with $Bk\mid(Ak)^2$.
The exhaustive counts are
\[
\begin{array}{c|rrrr}
t&15&30&45&60\\\hline
\theta\text{ totally positive}&0&6&28&69\\
\theta\text{ passing }d\text{-number test}&0&6&12&21\\
\text{admissible complementary pairs}&0&0&0&0.
\end{array}
\]
The verifier reconstructs the maximal order, all codegree embeddings and
ideal intersections by exact arithmetic; rational root intervals certify
all positivity decisions. This proves both general exclusions.

\paragraph{A common higher-field reduction for $7018$ and $8785$.}
For these two rings, respectively set
\[
 (c,T_0,B)=(195,4405,4),\qquad(65,4325,13).
\]
The original codegrees generate $E=\mathbb Q(\sqrt5)$ and their ideal
intersection is $c\sqrt5\mathcal O_E$. The full algebra modulo two is
$\mathbb F_4^3$: the Frobenius matrix squares to the identity and its
fixed space has dimension three. Therefore the signed order is etale
at two. Every completion of $K$ at two is unramified of degree two,
and its Frobenius restricts nontrivially to $E$. We have
\[
 [K:\mathbb Q]\in\{2,4,6\},\quad
 \Delta\in c\sqrt5\mathcal O_K\cap\mathfrak D_K,\quad
 \operatorname{Tr}_K(\Delta)\le T_0.
\]
The different condition is Theorem~\ref{thm:signed-different}:
orthogonality gives $\sum_i d_i\otimes d_i=(\Delta\otimes1)e_K$;
contracting against the inverse different shows
$\Delta\mathfrak D_K^{-1}\subset\mathcal O_K$.

A cyclic quartic $K$ is impossible because its order-two Frobenius
would act trivially on its unique quadratic subfield. In a biquadratic
field write $K=E\mathbb Q(\sqrt q)$, with $q>1$ odd squarefree,
$5\nmid q$ and $q\equiv1\pmod4$. The discriminants are coprime.
Writing
\[
 \Delta=c\sqrt5(a+b\sqrt5)/2,\qquad a,b\in\mathcal O_{\mathbb Q(\sqrt q)},
\]
shows that $b$ is totally positive with trace at most $B$.
If $p\mid q$ but $p\nmid c$, the different puts $b$ in the ramified
prime, so $p\mid\operatorname{Tr}(b)$. Thus
$q/\gcd(q,c)\mid\operatorname{Tr}(b)\le B$.
For $c=195$ the only possible prime factors of $q$ are $3,13$;
for $c=65$ they are $3,7,11,13$, with the product of those other than
$13$ at most $13$. The congruence $q\equiv1\pmod4$ leaves only $q=13$.

In this remaining field put
\[
 \Delta=\frac{c\sqrt5}{4}
 (x+y\sqrt5+z\sqrt{13}+w\sqrt{65}).
\]
The maximal-order conditions are that $x,y,z,w$ have the same parity
and $x-y-z+w\equiv0\pmod4$. Positivity gives the complete bounds
\[
 1\le y\le B,\quad x^2<5y^2,\quad13z^2<5y^2,\quad13w^2<y^2.
\]
Apply the degree-four $d$-number conditions and the complementary-pair
test above, now allowing quadratic radicands $5,13,65$.
There are seven $d$-number candidates for $7018$, with no pair.
For $8785$ there are thirty-five; the only two surviving pairs give
\[
 P_\Delta=(X^2-1625X+21125)^2,\qquad P_{f,g}=X^2-143X+169.
\]
Their product reduces to $(X^2+X+1)^3$, contradicting the required
original Casimir reduction $(X^2+X+1)(X+1)^4$.

For degree six, $K=EF$ with $F$ cyclic cubic and two completely split
in $F$. Its conductor has the complete form
$9^\epsilon\prod p$, where $\epsilon\in\{0,1\}$ and the distinct
primes satisfy $p\equiv1\pmod3$. An order-three character is trivial
on the pro-$p$ kernel for $p\ne3$ and factors through modulus nine at
three, which proves this form. The different exponents are $d_p=2$
at tame primes and $d_3=4$ at the block nine, by the
conductor--discriminant formula \cite[Chapter~V, Theorem~3.27]{MilneCFT}.
Again $b\in\mathcal O_F$ is totally positive with trace at most $B$.
The different forces $b\in\mathfrak p^{k_p}$, where
$k_p=\max(0,d_p-3v_p(c))$. The local codifferent definition gives
\[
 \operatorname{Tr}(\mathfrak p^k)
 =p^{\lfloor(k+d_p)/3\rfloor}\mathbb Z_p,
 \qquad
 \operatorname{Norm}(b)\ge\prod_{p\mid f}p^{k_p}.
\]
Indeed, containment of that trace ideal in $p^m\mathbb Z_p$ is
equivalent to $p^{-m}\mathfrak p^k\subset\mathfrak p^{-d_p}$.
For $7018$ these trace divisibilities leave only conductors $9,13,117$;
the first two make two inert. At $117$ the trace of $b$ is three and
its norm at least three, contradicting AM--GM's upper bound one.
For $8785$ the complete list is $7,9,13,91,117$: simultaneous blocks
$7,9$ would require trace divisible by $63$, and every other tame prime
is at least $19$. The single blocks make two inert. At $91$ the trace
is seven and norm at least $49>(7/3)^3$; at $117$ the trace is nine
and norm at least $81>27$. Thus both rings force $K=E$.

\paragraph{The quadratic spectra and the last signed search.}
For $7018$ write
\[
 \Delta=(975y+195x\sqrt5)/2,\qquad
 1\le y\le4,\quad x\equiv y\pmod2,\quad x^2<5y^2,
 \quad 5y^2-x^2\mid20y^2.
\]
There are sixteen possibilities. Because the etale two-adic algebra is
$\mathbb F_4^3$, every rational character component has even degree.
The four complementary codegrees therefore partition into two
$E$-conjugate pairs; rational pairs are repeated integers.
For a pair $f,f'=(S\pm b\sqrt5)/2$, set $P=ff'$ and require
\[
 P\mid\operatorname{Norm}(\Delta),\quad S^2-4P=5b^2,\quad
 S\equiv b\pmod2,\quad S>2,\quad P>S-1,
\]
together with $P\mid S^2$ when $b\ne0$, and integrality of
$\Delta/f,\Delta/f'$. The two pairs must have total trace at most
$4405-975y$ and reciprocal sum $1-\operatorname{Tr}_E(1/\Delta)$.
Only $(x,y)=(\pm1,1)$ remain, both with full polynomial
\[
 (X^2-975X+190125)(X^2-975X+1125)(X^2-975X+7605).
\]
Its reduction $(X^2+X+1)^3$ contradicts the original
$(X^2+X+1)(X+1)^4$, excluding $7018$.

For $8785$ the identical pair-enumeration procedure starts with
$\Delta=(325y+65x\sqrt5)/2$, $1\le y\le13$.
It leaves four spectra up to simultaneous quadratic conjugation.
Three have $\Delta=D=(4225+1885\sqrt5)/2$ or its conjugate, so
\cite[Proposition~8.23]{ENO} reduces them to the pseudo-unitary
obstruction below. The exceptional spectrum consists of the three pairs
\[
 \frac{1625\pm325\sqrt5}{2},\qquad
 \frac{450\pm200\sqrt5}{2},\qquad
 \frac{1170\pm520\sqrt5}{2}.
\]
Its signed Casimir trace is $3245$, hence its deficit is $1080$.
Lemma~\ref{lem:m9-indicator-gauge} leaves the explicit domain
\[
\begin{gathered}
 t_q\in\{-N_q,-N_q+2,\ldots,N_q\},\qquad
 \sum_qw_q(N_q^2-t_q^2)=1080,\\
 t_q\equiv N_q\pmod4\quad\text{if }q\text{ repeats an index}.
\end{gathered}
\]
There are $35$ nonunit triples. The independent verifier generates all
entries of all commutators of the five multiplication matrices, giving
$55$ distinct polynomial equations. Depth-first enumeration evaluates
each equation when its last variable is assigned and prunes only when
the nonnegative accumulated deficit exceeds $1080$. It visits
$2\,247\,571\,340$ nodes and finds no table. All domains, equations,
native integer bounds and successful termination are regenerated and
checked; this is a complete finite proof, not a search deadline.

\paragraph{Three rings with an involution.}
We first give their exact pseudo-unitary certificates in a common form.
For a commutative category let $G=\sum_iN_iN_{i^*}$ be the central
induction Gram matrix. The six summands of $I(1)$ have restriction
columns $v_j$ with unit coordinate one and dimension $D/f_j$.
Full-character reciprocity is
\[
 (D/f_i)\chi_i(v_j)=(D/f_j)\chi_j(v_i).
\]
For fixed low columns, the high columns and their relative twist moments
satisfy Proposition~\ref{prop:round3-full-character}. Its exact
nonnegative system $Ax=b$ is impossible whenever an integer vector
$y$ has $y^tA\ge0$ and $y^tb<0$.

For $529,2455$ a common complete low-column prune is as follows.
Let $u$ be the unit column, $M=G-uu^t>0$, fix one other low column $v$,
and put $w=G_{\bullet0}-u-v$. The four remaining columns sum to $w$,
so Cauchy's inequality gives
$M-vv^t-ww^t/4\succeq0$.
Writing $A=qM^{-1}$ integral with $q>0$, this is exactly
\[
 q-v^tAv\ge0,\quad4q-w^tAw\ge0,\quad
 (q-v^tAv)(4q-w^tAw)\ge(v^tAw)^2.
\]
The five column domains shrink from $(817,833,833,217,217)$ to
$(10,9,9,10,10)$. Exhausting their Cartesian product modulo
permutations of equal-codegree columns gives $497$ frames; ordinary
Gram/moment constraints reject $406$, and character reciprocity leaves
six frames for $529$ and thirty-one for $2455$.
Every one of the $8$, respectively $32$, possible second-indicator
assignments is retained. The resulting $48$ and $992$ systems have
exact integer Farkas witnesses.

For $1060$, use the basis $(1,g,X,Y,gY,gX)$.
The five nonunit low columns have total multiplicities in the paired
$X,gX$ and $Y,gY$ coordinates
\[
 A=(16,11,11,6,22),\qquad B=(17,11,11,5,22).
\]
Their difference vectors $u,v$ obey
$\sum u_j=\sum v_j=-2$, $\sum u_j^2\le12$,
$\sum v_j^2\le20$, $u\equiv A$, $v\equiv B\pmod2$.
There are $18\cdot59=1062$ pairs; positivity of the remaining Gram
matrix leaves $369$, and exact character reciprocity leaves four
frames. All $32$ indicator choices for each frame are contradicted by
$128$ integral Farkas witnesses. These three proofs use no indicator
parity shortcut and no numerical optimization in replay.

For the general reduction of $1060$, positive sphericalization lifts
make $g$ act positively on its three free pairs. The dimension character
is $(1,1,x,y,y,x)$, so its field is a quotient of the rank-three
$g=1$ algebra. It contains $\mathbb Q(\sqrt{273})$, hence equals that
field. Write $x=(a+b\sqrt{273})/2$, $y=(c+d\sqrt{273})/2$.
Orthogonality with the conjugate character gives
\[
 a^2+c^2+4=273(b^2+d^2),\qquad
 \Delta=273s+(ab+cd)\sqrt{273},\quad s=b^2+d^2.
\]
The original trace is $4396$ and the other four codegrees exceed one,
so $0<s\le8$. Thus $b,d\in[-2,2]$, $a,c\in[-46,46]$.
There are $144$ parity-compatible solutions. Integrality of
$\Delta/12$ and $\Delta/(2184+132\sqrt{273})$ leaves $32$ assignments
and the four values $2184\pm108\sqrt{273}$,
$2184\pm132\sqrt{273}$. The first pair fails the quadratic
$d$-number test: its norm $1585584$ does not divide $4368^2$.
The second pair is the FP global-dimension orbit. Galois conjugation
and \cite[Proposition~8.23]{ENO} contradict the pseudo-unitary
certificate, proving general nonexistence.

For $529,2455$, use $(1,g,x,y,a,b)$ with $gx=y$, $ga=a$, $gb=b$.
The FP dimensions are $(1,1,d,d,d-1,d+1)$, $d=16+\sqrt{257}$,
and the codegrees are $4,4,8,8,D,D'$ with
$D=2056+128\sqrt{257}$. The original trace is $4136$.
Gauge the lift of $g$ to dimension one; its fixed objects force its
left and right action signs to agree. In $529$, $x^*=y$ and dual
invariance forces positive action on that pair; in $2455$ gauge the
pair to achieve it. The signed $g=1$ and $g=-1$ components have
dimensions four and two. The former has a one-dimensional character,
so neither component can contain a noncommutative matrix block.
Thus the signed algebra is commutative, with $E=\mathbb Q(\sqrt{257})
\subset K$ and $[K:\mathbb Q]=2$ or $4$.
For $\theta=\Delta/8$, induction gives
$\theta,\theta/\sqrt{257}\in\mathcal O_K$.

If $K=E$, the trace is a positive multiple of $2056$.
At both split two-adic embeddings, the mod-two table gives $x$ odd and
$a=2a_0,b=2b_0$. Integrality of
$\Delta/8=(2+2x^2+4(a_0^2+b_0^2))/8$ forces
$a_0^2+b_0^2$ odd, hence one modulo four. Therefore
$\Delta\equiv8\pmod{16}$ at both embeddings, and its trace is
divisible by sixteen. The only possible trace is $4112$.
The four other codegrees are positive. Decreasing any coefficient
magnitude $7,8,9$ loses at least $24$ from the full trace, which
would leave at most $4112$. Thus the signed table is saturated.
The complete oriented sign domains, with positive $g$ action, contain
exactly eight associative tables for each ring; every table has a
certified unit- and dual-preserving gauge to the original one.
Each therefore supplies a spherical section. The nonzero dimension
characters lie only in the FP orbit (the others have zero entries),
so the pseudo-unitary certificates exclude this case.

If $[K:\mathbb Q]=4$, the canonical trace of $(1+g)/2$ gives
$\operatorname{Tr}_K(1/\theta)=4$ and
$\operatorname{Tr}_K(\theta)\le517$.
The element $\theta$ generates $K$: if quadratic of norm $n$, these
identities give full trace $4n$, whereas
$\theta/\sqrt{257}$ integral forces $n\ge257$; if rational, they
give $\theta=1$, also impossible. Relative trace and integrality give
the minimal polynomial
\[
 P=X^4-257uX^3+257wX^2-4\cdot257^2vX+257^2v,
\]
where $u=1,2$ and
\[
 v\mid257^2u^4,\quad v^2\mid w^4,\quad
 1\le v\le\lfloor257u^3/64\rfloor,\quad
 1\le w\le\lfloor3\cdot257u^2/8\rfloor.
\]
The bounds are Maclaurin's inequalities, and the divisibilities are
the $d$-number test. Exactly $1057$ polynomials occur; exact Sturm
counts reject $736$ for lack of four positive roots. The discriminants
of the remaining $321$ are neither squares nor $257$ times squares.
These are the two possibilities for an abelian quartic containing $E$
(biquadratic or cyclic). This excludes the quartic case.

\paragraph{The pseudo-unitary obstruction for $8785$.}
Put $\phi=(1+\sqrt5)/2$, $t=8+13\phi$ and
$D=5(1+t^2)$. The dimensions are
$(1,t,t-1,t-1,t+1,t+1)$ and the two smallest nonunit induction
dimensions are both $3t+2$. Their columns have unit coordinate one,
the other five coordinates summing to three, and exactly eight
nonnegative solutions to the dimension equation. For each of the
$8^2$ labeled pairs, dimensions and reciprocity with these two
columns give a rational system of rank six for every other column.
Exactly ten nonnegative integral frames remain; full row sums and
reciprocity are checked.

For a nontrivial self-dual simple of dimension $d$, let $\ell$ be the
dimension mass in its induction contributed by the twist-one summands
of $I(1)$. By \cite[Theorems~2.5 and~2.7]{Ostrik3}, the first twist
moment is zero and the second is $D\nu_2(X)$, with $\nu_2(X)\le1$.
Cauchy's inequality applied to the remaining real twist moments gives
\[
 2\ell^2\le(Dd-\ell)\bigl(D(d+1)-2\ell\bigr),\qquad
 (3d+1)\ell\le Dd(d+1).
\]
For $d=t$, the ten frames have low masses
$14473+23413\phi$ (six frames), $13641+22061\phi$ (two), or
$12809+20709\phi$ (two). The differences between the left and right
sides of the last inequality are respectively
\[
 257812+417144\phi,\quad184284+298168\phi,\quad110756+179192\phi,
\]
all positive. This excludes every frame and completes the $8785$ proof.

\paragraph{The rank-three extension $3748$: categorical reduction.}
This tensor is
\[
 R=A[X]/(X^2-1-QX),\qquad Q=2+2a+b,
 \quad a^2=1+b,\ ab=a+b,\ b^2=1+a+b.
\]
The basis $(1,a,b,X,aX,bX)$ is source order $(0,1,5,2,3,4)$.
In a signed lift of $A$, write
$a^2=1+ub$, $ab=ua+vb$, $b^2=1+va+wb$, with $u,v,w=\pm1$.
Associativity gives $uw=1$, and a sign gauge makes all signs positive.
Thus any categorification of $A$ is spherical, and
\cite[Theorem~1.1]{Ostrik3} identifies its fusion category
$\mathcal D$ as $\operatorname{PSU}(2)_5$ or a Galois conjugate.

Lift this fusion-closed $A$ to sphericalization of the full category,
choose $X$, and define the other lifts by $aX,bX$. Its entire $A$ action
is positive. The signed relation is
$X^2=1+(p+qa+rb)X$, where $p,q\in\{-2,0,2\}$ and $r=\pm1$.
Simultaneously changing the three complementary lifts allows $r=1$.
For the nine cases the signed Casimir polynomials are irreducible
sextics. Their pairs (trace, constant coefficient) are
\[
\begin{gathered}
 (189,3296573),(77,703493),(245,6754013),\\
 (105,1279733),(105,434581),(385,2854789),\\
 (133,1279733),(245,3027661),(637,5274997).
\end{gathered}
\]
In lexicographic order on $(p,q)$, the first eight violate
$P(0)\mid\operatorname{Tr}(P)^6$. Thus the sphericalization
$d$-number test leaves only $(p,q,r)=(2,2,1)$, giving a spherical
section of the full category. Its original global-codegree polynomial
\[
 P_f=X^6-637X^5+22295X^4-298116X^3
        +1854944X^2-5274997X+5274997
\]
is irreducible, so there is a pseudo-unitary Galois conjugate.

Let $\mathcal M$ be the complementary $\mathcal D$-bimodule category.
The actual left-module functor $d\mapsto dX$ maps simples bijectively
and is an equivalence. Transport of its right action gives a tensor
autoequivalence of $\mathcal D$ fixing simple labels; hence
$\mathcal M$ is an invertible twisted regular bimodule.
The result $\operatorname{BrPic}(\operatorname{Ad}(A_6))=1$,
including Galois conjugates, is \cite[Theorem~1.2]{EMADE}.
Thus $\mathcal M$ is categorically regular, with the unit sent to $X$.

Its relative center therefore has simple basis $B\sqcup Bt$, where
$B=A\otimes A=K_0(Z(\mathcal D))$ and $t$ forgets to $X$.
The central action of $Z(\mathcal D)$, equivalently naturality of the
relative half-braiding, makes $B$ commute with $t$.
The unique simple over $X$ is self-dual, and positivity forces
\[
 t^2=1+Ct,\qquad C\in B_{\ge0},\qquad F(C)=2+2a+b.
\]
These are the fifteen explicitly enumerated relative-center lifts.
The formal-codegree $d$-number test excludes thirteen, leaving only
$C=Q\otimes1$ and $C=1\otimes Q$.
Each one-sided case contains a fusion subcategory whose forgetful
functor to the full category is bijective on simples, hence a tensor
equivalence. Its section supplies coherent half-braidings with all
objects, giving a central lift of $\mathcal D$ to $Z(\mathcal C)$;
on $\mathcal D$ it is one of the two nondegenerate braidings.
The relative-center construction is \cite[Section~2.2]{GNNrel}.
Consequently \cite[Theorem~3.13]{DGNO} applies inside the center:
\[
 Z(\mathcal C)=\mathcal D\boxtimes\mathcal E.
\]
No braiding or tensor factorization of $\mathcal C$ has been assumed.

\paragraph{A six-versus-four Galois contradiction.}
Use positive dimensions
\[
 a=2\cos(\pi/7),\quad b=a^2-1,\quad h=\sqrt{13},\quad
 q=2+2a+b,\quad x=(q+bh)/2,
\]
and put $d=1+a^2+b^2$, $f=d(1+x^2)$.
The identities $a^3-a^2-2a+1=0$, $q^2+4=13b^2$ and
$x^2-qx-1=0$ hold. The six numbers $1,a,b,x,ax,bx$ are
$\mathbb Q$-independent, so dimensions determine forgetful vectors.
The centralizer has dimension $f^2/d$.

The induction object $I(1)$ has six multiplicity-one simple summands,
all of twist one, by \cite[Theorems~2.5 and~2.13]{Ostrik3}.
Its integral multiplicity vector $n$ satisfies normalized $Sn=n$
by \cite[Proposition~2.8]{Ostrik3}, so each signed Galois permutation
fixes $n$. The six distinct unit Galois dimensions $f/\sigma(f)$,
distinct because $P_f$ is irreducible, exhaust these summands.
Normalized $S$ is compatible with the Deligne product. The three
embeddings $a\mapsto a,a^2-a-1,2-a^2$ have $\mathcal D$ unit
Galois labels $1,b,a$, respectively, since their squared dimensions
are $d/\sigma(d)$. Dividing by these dimensions identifies six known
$\mathcal E$ simples with forgetful vectors
\[
\begin{gathered}
 (1,0,0,0,0,0),\quad(1,0,0,2,2,1),\\
 (0,0,1,1,3,4),\quad(0,0,1,0,0,1),\\
 (0,1,0,1,2,1),\quad(0,1,0,1,1,2).
\end{gathered}
\]
Their twists are inverse to those of their $\mathcal D$ labels.
Since $\operatorname{Norm}(d)=49$, Cauchy's theorem
\cite[Theorem~3.9]{BNRW} makes all these twists seven-primary.
The second object, $e_-$, corresponds to $h\mapsto-h$, and has
dimension $x^2$ and twist one.

The induction Gram identity is
$\sum_{z\in\operatorname{Irr}Z(\mathcal C)}[Fz][Fz]^t=\sum_iN_i^2$.
Subtract the eighteen vectors obtained by multiplying these six known
vectors by $1,a,b$. The residual is zero in the first three rows and
columns and has lower block
\[
 T=\begin{pmatrix}24&36&44\\36&68&80\\44&80&104\end{pmatrix}.
\]
All contributions are nonnegative outer products. Every unknown
$\mathcal E$ simple therefore forgets to $XU$, where
$U=u_0+u_1a+u_2b$, $u_i\ge0$, not all zero.
Its whole $\mathcal D$ block contributes $u_0^2+u_1^2+u_2^2$ to
$T_{00}$, so the sum of these squared norms over unknown simples is
$24$.

Under $a\mapsto a^2-a-1$, $h\mapsto h$, the positive dimension
Galois action on unknown vectors is $u\mapsto\pm Hu$, where
\[
 H=\begin{pmatrix}0&-1&1\\-1&-1&1\\-2&0&1\end{pmatrix},\qquad H^3=-I.
\]
This follows directly from
$d_{\widehat\sigma(y)}=d_{\widehat\sigma(1)}|\sigma(d_y)|$:
the transformed $U$ is $(b^2/\sigma(b))\sigma(U)$.
Injectivity of the dimension map requires all coordinates to have
one sign; mixed signs are impossible. A dimension orbit has length one
or three, and equal multiplicity along the orbit.
A fixed positive vector is $n(1,2,1)$, with squared norm $6n^2$.
If a three-element orbit had norm sum below six, one vector would be
a coordinate unit. Their complete positive-sign orbits have norm sums
$12,12,6$, respectively. Hence every dimension orbit costs at least
six in the Gram budget. Any exact dimension can therefore occur with
multiplicity at most $24/6=4$.

On the other hand,
$\operatorname{Norm}_{\mathbb Q(a,h)/\mathbb Q}(f^2/d)=7^4\,13^6$.
Cauchy's theorem forces an $\mathcal E$ twist with a nontrivial
thirteen-primary component; the corresponding object is unknown.
Choose a cyclotomic automorphism fixing all prime-to-thirteen roots
and acting by exponent two on the thirteen-primary roots, possible by
the Chinese remainder theorem in a common cyclotomic field for the
normalized modular representation and its twists.
It fixes $a$, negates $h$, and sends $x$ to $-1/x$.
Its unit image is $e_-$, so it preserves each unknown dimension:
$x^2|\sigma(xU)|=xU$.
For actual twists the normalized modular Galois theorem
\cite[Theorem~II(iii)]{DLN}, divided by its unit case, reads
\[
 \theta_{\widehat\sigma(y)}
 =\theta_{\widehat\sigma(1)}\sigma^2(\theta_y)
 =\sigma^2(\theta_y).
\]
Exponent four has order six modulo thirteen. Thus the nontrivial
thirteen-primary twist forces at least six different objects of the
same dimension, also for higher thirteen powers. This contradicts
the bound four and completes the proof of
Theorem~\ref{thm:m9-completion}.
